\documentclass[14pt]{beamer}
\usepackage{pepslides}
\title{Delay time tomography}
\lecturenumber{19}
\hypersetup{pdftitle={Lecture 19: Delay time tomography}}

\begin{document}

\titleslide

\begin{frame}[t]
\heading{From last time\ldots}
\centering
\includegraphics[height=0.7\paperheight]{seismology_pioneers.jpg}
\end{frame}

\sectionslide{Lateral velocity variations}

\begin{frame}
\centering
\includegraphics[height=0.85\paperheight]{mantle_convection_cartoon.jpg}
\end{frame}

\begin{frame}
\centering
\includegraphics[height=0.75\paperheight]{mantle_convection_models.jpg}

Can we do better than drawing cartoons?
\end{frame}

\begin{frame}
\begin{columns}[c]
\begin{column}{0.48\textwidth}
\centering
\includegraphics[width=\linewidth]{L7F1.png}
\end{column}
\begin{column}{0.48\textwidth}
\textbf{Seismic tomography}: quantitative models for lateral variations within the Earth derived from seismological
observations.

Most models show perturbations in elastic wave speeds relative to a spherically symmetric reference model.

Other parameters (e.g.\ density, attenuation) can be sought with more difficulty.

Simple interpretation:

\textbf{Red = slow = hot}

\textbf{Blue = fast = cold}
\end{column}
\end{columns}
\credit{Woodhouse \& Dziewonski (1989)}
\end{frame}

\sectionslide{Setting up the inverse problem}

\begin{frame}
The \textbf{delay time} is the measured travel time minus that predicted by a spherical reference model:
\begin{center}
\includegraphics[width=0.75\textwidth]{delay_time_seismograms.jpg}
\end{center}
\begin{eqn}
\delta T = T - T_{0}.
\end{eqn}
\end{frame}

\begin{frame}
Rays determined through integration of Hamilton's equations:
\begin{eqn}
\frac{\ddns x_{i}}{\ddns\sigma}=\frac{\partial H_{k}}{\partial p_{i}}, \quad \frac{\ddns p_{i}}{\ddns\sigma}=-\frac{\partial H_{k}}{\partial x_{i}}.
\end{eqn}
Once the ray path is known, the travel time can be written:
\begin{eqn}
T = \int_{\mathrm{ray}} \frac{1}{v} \dd s.
\end{eqn}
If velocity variations are sufficiently small, we can seek to linearise:
\begin{eqn}
T(v+\delta v) = T(v) + \langle DT(v)|\delta v \rangle + O(\delta v^{2}).
\end{eqn}
\end{frame}

\begin{frame}
Using Fermat's principle, we find that:
\begin{eqn}
\delta T = -\int_{\mathrm{ray}} \frac{\delta v}{v^{2}} \dd s.
\end{eqn}
First-order contribution from varying ray geometry vanishes, and so we need only
integrate the velocity perturbation along the reference ray path.

Functional derivative identified by using the standard inner product:
\begin{eqn}
\langle u|w \rangle = \int_{M} uw \dd^{3}\mathbf{x}.
\end{eqn}
Note that the functional derivative is singular on the ray path and zero elsewhere.
\end{frame}

\begin{frame}
For the $i$th delay time measurement:
\begin{eqn}
\delta T_{i} = -\int_{\mathrm{ray}_{i}} \frac{\delta\alpha}{\alpha^{2}} \dd s + e_{i}.
\end{eqn}
Expand the velocity perturbation in a finite-dimensional basis:
\begin{eqn}
\delta\alpha(\mathbf{x}) = \sum_{j=1}^{m} \delta\alpha_{j} \varphi_{j}(\mathbf{x}).
\end{eqn}
Choice of basis functions determined by assumptions about the Earth and the potential
resolution of the data as set by its dominant frequencies.
\end{frame}

\begin{frame}
Using the model parametrisation, we now have:
\begin{eqn}
\delta T_{i} = \sum_{j=1}^{m} \left(-\int_{\mathrm{ray}_{i}} \frac{\varphi_{j}}{\alpha^{2}} \dd s \right) \delta\alpha_{j} + e_{i}.
\end{eqn}
We can now use a familiar matrix-vector notation:
\begin{eqn}
\mathbf{d} = \mathbf{A}\mathbf{m} + \mathbf{e},
\end{eqn}
where we have set:
\begin{eqn}
A_{ij} = -\int_{\mathrm{ray}_{i}} \frac{\varphi_{j}}{\alpha^{2}} \dd s.
\end{eqn}
\end{frame}

\begin{frame}
A toy problem: straight rays crossing a square divided into $24\times24$ blocks, a \textbf{pixel basis}.
\begin{center}
\includegraphics[width=0.85\textwidth]{L7F2.pdf}
\end{center}
$A_{ij}$ is minus the length of the $i$th ray within the $j$th block divided by $\alpha^{2}$, and so most entries are zero.
\end{frame}

\sectionslide{Solving the inverse problem}

\begin{frame}
Tomographic problem with perfect data:
\begin{eqn}
\mathbf{d} = \mathbf{A}\mathbf{m}.
\end{eqn}
Is there a unique model that solves this problem?

The \textbf{null space} of the matrix $\mathbf{A}$ comprises those vectors satisfying:
\begin{eqn}
\mathbf{A}\mathbf{m}_{\mathrm{null}} = \mathbf{0}.
\end{eqn}
This is a linear subspace. If non-trivial (not just the zero vector), then the linear system
cannot have a unique solution.

\textbf{The null space is always non-trivial for under-determined problems, $\bm{m > n}$.}
\end{frame}

\begin{frame}
\centering
\includegraphics[width=0.9\textwidth]{L7F3.pdf}

\raggedright
Lack of ray coverage is the most obvious way null spaces are generated.

Adding the null-space model (left) to the input model gives exactly the same data.
\end{frame}

\begin{frame}
The null space can be found from the eigenvalue problem for the symmetric matrix $\mathbf{A}^{T}\mathbf{A}$, whose
eigenvalues are non-negative:
\begin{eqn}
\mathbf{A}^{T}\mathbf{A}\mathbf{v}_{j}=s_{j}^{2}\mathbf{v}_{j}.
\end{eqn}
The null space is spanned by the eigenvectors with $s_{j}=0$.

Setting $\mathbf{u}_{j}=\mathbf{A}\mathbf{v}_{j}/s_{j}$ for $s_{j}>0$, we obtain the \textbf{singular value decomposition}:
\begin{eqn}
\mathbf{A}=\sum_{s_{j}>0}s_{j}\mathbf{u}_{j}\mathbf{v}_{j}^{T}.
\end{eqn}
The data depend only on the components of the model along $\mathbf{v}_{j}$ with $s_{j}>0$.
\end{frame}

\begin{frame}
\centering
\includegraphics[width=\textwidth]{L7F4.pdf}

\raggedright
Directions with small singular values contribute little to the data, and are lost within the errors: they are
effectively unresolved.
\end{frame}

\begin{frame}
An obvious way to proceed is through a \textbf{least-squares solution}.

We define the least-squares misfit through:
\begin{eqn}
J(\mathbf{m}) = \frac{1}{2} \sum_{i=1}^{n} \frac{[(\mathbf{A}\mathbf{m})_{i} - d_{i}]^{2}}{\sigma_{i}^{2}}.
\end{eqn}
In matrix notation, this becomes:
\begin{eqn}
J(\mathbf{m}) = \frac{1}{2} (\mathbf{A}\mathbf{m} - \mathbf{d})^{T} \mathbf{C}^{-1} (\mathbf{A}\mathbf{m} - \mathbf{d}).
\end{eqn}
If there is a non-trivial null space there is \textbf{no unique minimum}!
\end{frame}

\begin{frame}
The idea is to penalise models with undesirable properties within the inversion.

Suppose we prefer models with a small norm:
\begin{eqn}
\int_{M} \delta\alpha^{2} \dd^{3}\mathbf{x}.
\end{eqn}
Relative to our model parametrisation, this becomes:
\begin{eqn}
\begin{gathered}
\int_{M} \delta\alpha^{2} \dd^{3}\mathbf{x} = \mathbf{m}^{T}\mathbf{B}\mathbf{m},\\[\eqgap]
B_{jk} = \int_{M} \varphi_{j}\varphi_{k} \dd^{3}\mathbf{x}.
\end{gathered}
\end{eqn}
\end{frame}

\begin{frame}
We modify the least-squares misfit to incorporate regularisation:
\begin{eqn}
J(\mathbf{m}) = \frac{1}{2} (\mathbf{A}\mathbf{m} - \mathbf{d})^{T} \mathbf{C}^{-1} (\mathbf{A}\mathbf{m} -
\mathbf{d}) + \frac{\lambda}{2} \mathbf{m}^{T}\mathbf{B}\mathbf{m}.
\end{eqn}
Models with a small value of $J$ now represent a balance between fitting the data well and
keeping the chosen model property small.

This balance is reflected in the choice of the \textbf{damping parameter}, $\lambda$.

Many choices of regularisation possible such as \textbf{norm or gradient damping}.
\end{frame}

\begin{frame}
\begin{columns}[c]
\begin{column}{0.5\textwidth}
\centering
\includegraphics[width=\linewidth]{L7F5.pdf}
\end{column}
\begin{column}{0.46\textwidth}
How should the damping parameter be chosen?

Fit the data to within their errors, but no better:
\begin{eqn}
\chi^{2}\approx n.
\end{eqn}
This is the \textbf{discrepancy principle}.

Fitting the data more closely than this serves only to fit the noise.
\end{column}
\end{columns}
\end{frame}

\begin{frame}
Misfit is quadratic, and so has at most one minimum value when:
\begin{eqn}
\nabla J(\mathbf{m}) = \mathbf{0}.
\end{eqn}
Through a calculation (detailed in the handout) this implies:
\begin{eqn}
(\mathbf{A}^{T}\mathbf{C}^{-1}\mathbf{A} + \lambda \mathbf{B})\mathbf{m} = \mathbf{A}^{T}\mathbf{C}^{-1}\mathbf{d}.
\end{eqn}
If the matrix on the left-hand side is invertible, there is a unique solution to the problem.
\end{frame}

\begin{frame}
Consider the associated eigenvalue problem:
\begin{eqn}
(\mathbf{A}^{T}\mathbf{C}^{-1}\mathbf{A} + \lambda \mathbf{B})\mathbf{m} = \mu \mathbf{m}.
\end{eqn}
If the eigenvalues are all non-zero, then the matrix is invertible.

Take the inner product with a normalised eigenvector:
\begin{eqn}
(\mathbf{A}\mathbf{m})^{T}\mathbf{C}^{-1}(\mathbf{A}\mathbf{m}) + \lambda \mathbf{m}^{T}\mathbf{B}\mathbf{m} = \mu.
\end{eqn}
The left-hand side is positive so long as the damping parameter is positive.

Regularised least squares is a means of calculating solutions to the inverse problem that
also have desirable physical properties. But there are a \textbf{lot of arbitrary choices}.
\end{frame}

\begin{frame}
Taking $\mathbf{B}=\mathbf{I}$ and $\mathbf{C}=\sigma^{2}\mathbf{I}$, the solution can be written using the singular value decomposition:
\begin{eqn}
\mathbf{m}=\sum_{s_{j}>0}\frac{s_{j}}{s_{j}^{2}+\lambda\sigma^{2}}(\mathbf{u}_{j}\cdot\mathbf{d})\,\mathbf{v}_{j}.
\end{eqn}
Components with $s_{j}^{2}\gg\lambda\sigma^{2}$ are recovered, while those with $s_{j}^{2}\ll\lambda\sigma^{2}$ are suppressed.

Damping acts by discarding those parts of the model to which the data are least sensitive.
\end{frame}

\begin{frame}
We can also ask that the model stays ``close'' to a preferred point:
\begin{eqn}
J(\mathbf{m}) = \frac{1}{2}(\mathbf{A}\mathbf{m}-\mathbf{d})^{T}\mathbf{C}^{-1}(\mathbf{A}\mathbf{m}-\mathbf{d})
+ \frac{\lambda}{2}(\mathbf{m}-\mathbf{m}_{0})^{T}\mathbf{B}(\mathbf{m}-\mathbf{m}_{0}).
\end{eqn}
For positive damping parameter, we obtain the unique solution:
\begin{eqn}
\mathbf{m} = (\mathbf{A}^{T}\mathbf{C}^{-1}\mathbf{A} + \lambda \mathbf{B})^{-1} (\mathbf{A}^{T}\mathbf{C}^{-1}\mathbf{d}
+ \lambda \mathbf{B}\mathbf{m}_{0}).
\end{eqn}
\end{frame}

\begin{frame}
\begin{columns}[c]
\begin{column}{0.5\textwidth}
\centering
\includegraphics[width=\linewidth]{L7F6.pdf}
\end{column}
\begin{column}{0.46\textwidth}
A \textbf{spike test} inverts synthetic data for a model that is zero except within a single block.

The recovered spike is reduced in amplitude and smeared along the rays, to an extent set by the ray coverage.
\end{column}
\end{columns}
\end{frame}

\begin{frame}
\centering
\includegraphics[height=0.74\paperheight]{dewit2012_nullspace_model.jpg}

Which features in the model can be believed? Not an easy question!
\credit{de Wit \emph{et al.} (2012)}
\end{frame}

\begin{frame}
\centering
\includegraphics[height=0.74\paperheight]{dewit2012_nullspace_robust.jpg}

Can vary arbitrary parameters to see what has to be there to fit the data.
\credit{de Wit \emph{et al.} (2012)}
\end{frame}

\begin{frame}
\centering
\includegraphics[height=0.74\paperheight]{L12F5.png}

Different data sets and methods lead to broadly consistent results.
\credit{Ritsema \emph{et al.} (2011)}
\end{frame}

\sectionslide{A Bayesian perspective}

\begin{frame}
\textbf{Bayes' theorem} expressed in terms of conditional probability density functions:
\begin{eqn}
p(\mathbf{m}|\mathbf{d}) = \frac{ p(\mathbf{d}|\mathbf{m}) p(\mathbf{m})}{\int p(\mathbf{d}|\mathbf{m}) p(\mathbf{m}) \dd \mathbf{m}}.
\end{eqn}
This is a complete recipe to update prior beliefs given data.

But what do probabilities mean in this context?

And how do you arrive at the ``correct prior''?
\end{frame}

\begin{frame}
\centering
\begin{minipage}{0.72\textwidth}
\emph{Now any decision that depends on the data that is being used in making the inference only requires from the data
the posterior distribution. Consequently the problem of inference is effectively solved by stating the posteriori
distribution. This is the reason I feel that the basic problem of inference is solved.}

D.~V. Lindley, \emph{Foundations of Statistical Inference}

\vspace{0.5em}
\emph{I am often asked if the method gives the right answer; or, more particularly, how do you know if you have got the
right prior. My reply is that I don't know what is meant by ``right'' in this context. The Bayesian theory is about
coherence, not about right or wrong.}

D.~V. Lindley, \emph{Bayesian Statistics}
\end{minipage}
\end{frame}

\begin{frame}
Prior and likelihood take the form:
\begin{eqn}
\begin{gathered}
p(\mathbf{m}) = \frac{1}{\sqrt{(2\pi)^{m}\det \mathbf{C}_{m}}} \exp\left[-\frac{1}{2}(\mathbf{m}-\mathbf{m}_{0})^{T}
\mathbf{C}_{m}^{-1}(\mathbf{m}-\mathbf{m}_{0})\right],\\[\eqgap]
p(\mathbf{d}|\mathbf{m}) = \frac{1}{\sqrt{(2\pi)^{n}\det \mathbf{C}}} \exp\left[-\frac{1}{2}(\mathbf{A}\mathbf{m}
-\mathbf{d})^{T}\mathbf{C}^{-1}(\mathbf{A}\mathbf{m}-\mathbf{d})\right].
\end{gathered}
\end{eqn}
Bayes' theorem results in a posterior PDF that is also Gaussian.
\end{frame}

\begin{frame}
Posterior mean has the same form as a regularised least-squares solution:
\begin{eqn}
\mathbf{m}_{p} = (\mathbf{A}^{T}\mathbf{C}^{-1}\mathbf{A} + \mathbf{C}_{m}^{-1})^{-1} (\mathbf{A}^{T}
\mathbf{C}^{-1}\mathbf{d} + \mathbf{C}_{m}^{-1}\mathbf{m}_{0}).
\end{eqn}
Method also provides the posterior covariance:
\begin{eqn}
\mathbf{C}_{p} = (\mathbf{A}^{T}\mathbf{C}^{-1}\mathbf{A} + \mathbf{C}_{m}^{-1})^{-1}.
\end{eqn}
Crucially, this quantifies uncertainty on the model -- a combination of propagated data
errors and the limitations imposed by the physics of the problem.

Any choice of regularisation amounts to a choice of prior: the Bayesian approach makes this choice explicit.
\end{frame}

\begin{frame}
Toy problem with a smooth Gaussian prior:
\begin{center}
\includegraphics[width=\textwidth]{L7F7.pdf}
\end{center}
The posterior standard deviation shows where the recovery can be trusted.
\end{frame}

\begin{frame}
Alternatively, pose and solve the problem for the function $\delta\alpha$ itself, discretising only at the end:
\begin{center}
\includegraphics[width=\textwidth]{L7F8.pdf}
\end{center}
You need only the idea: the answer no longer depends on an arbitrary parametrisation.
\credit{Calculated with \texttt{pygeoinf}}
\end{frame}

\begin{frame}[t]
\heading{What you need to know and be able to do}
\begin{enumerate}
\item[(i)] Why the construction of tomographic models is of interest: it gives the main quantitative constraint on
mantle convection and hence on the dynamics of the Earth's deep interior. Here you should be familiar with the idea
that seismic slow regions are likely to be hotter than average, while fast regions will be cooler.
\item[(ii)] How Fermat's principle is used to obtain a simple linearised relation between delay times and velocity
perturbations. From this, you should be able to outline how a tomographic problem is reduced to the standard form
$\mathbf{d} = \mathbf{A}\mathbf{m} + \mathbf{e}$.
\item[(iii)] What the null space of a linear(ised) inverse problem is and why it matters.
\item[(iv)] How regularised least squares is used to obtain solutions of linearised tomographic problems. Here you
should be able to derive the least squares equations from first principles.
\item[(v)] How Bayes' theorem can be applied in the solution of inverse problems, and that the results obtained
coincide with regularised least squares for linear(ised) problems with Gaussian errors and priors.
\end{enumerate}
\end{frame}

\end{document}
